\chapter{Free-boundary geometry in the ball}\label{sec:free-boundary-ball}
These local calculations connect the first variation, boundary spectra, and instability. The geometric source is Fraser's free-boundary variational paper \cite{Fraser00}, Section 2, especially Theorem 2.9, pp.~967--970. The spectral viewpoint is \cite{FS11}, Sections 1 and 3.

\begin{definition}[Free boundary in a Euclidean ball]\label{def:ball-free-boundary}
Let $F:\Sigma^n\to\overline B^N$ be a smooth immersion of a compact manifold with nonempty boundary, where $1\leq n<N$. We require $F(\operatorname{int}\Sigma)\subset B^N$, $F(\partial\Sigma)\subset S^{N-1}$, and that its outward unit conormal, viewed in the ambient Euclidean space, satisfies $\eta=F$ on $\partial\Sigma$. The immersion is a free-boundary minimal immersion if its mean curvature vector vanishes. The conormal condition says that the surface meets the supporting sphere orthogonally. All volumes below use the induced metric, including multiplicity for an immersion.
\end{definition}

\begin{proposition}[Stationarity for boundary-preserving variations]\label{prop:ball-free-boundary-stationarity}
A free-boundary minimal immersion is stationary under every smooth variation whose boundary remains on the unit sphere. In particular, this holds for ambient flows generated by vector fields tangent to the sphere.
\end{proposition}
\begin{proof}
If $V=\partial_tF_t|_0$ is the variation field, differentiating the induced metric gives $A'(0)=\int_\Sigma\operatorname{div}_\Sigma V$. Decompose into tangential and normal parts and apply the divergence theorem:
\[
 A'(0)=-\int_\Sigma\langle\vec H,V\rangle+\int_{\partial\Sigma}\langle\eta,V\rangle.
\]
The interior term is zero by minimality. Differentiating $|F_t|^2=1$ on the boundary gives $\langle F,V\rangle=0$ there, and $\eta=F$ kills the boundary term. No fixed-boundary condition is imposed on $V$.
\end{proof}

\begin{proposition}[Coordinate functions and the area identity]\label{prop:free-boundary-coordinate-identities}
For a free-boundary minimal immersion into the unit ball,
\[
 \Delta_\Sigma F=0,\qquad \partial_\eta F=F,\qquad
 \int_{\partial\Sigma}F=0,\qquad
 |\partial\Sigma|=n|\Sigma|.
\]
The vector boundary integral is taken componentwise.
\end{proposition}
\begin{proof}
Tracing the covariant Hessian of an immersion gives $\Delta F=\vec H=0$. Since $dF(\eta)=F$ along the boundary, its coordinate functions have the displayed Robin condition. Integrating $\Delta F=0$ gives the zero boundary integral. Also $|dF|^2=n$, hence $\Delta|F|^2=2n$. Its outward derivative is $2\langle F,dF(\eta)\rangle=2$ on the boundary. Integration proves the area identity. These formulas use the trace convention for $\vec H$, not its average.
\end{proof}

\begin{definition}[The Steklov equation]\label{def:steklov-equation}
On a compact connected Riemannian manifold $\Sigma$ with nonempty boundary, a nonzero smooth real function $u$ is a Steklov eigenfunction with eigenvalue $\sigma$ if
\[
 \Delta u=0\quad\hbox{in }\Sigma,\qquad
 \partial_\eta u=\sigma u\quad\hbox{on }\partial\Sigma.
\]
This section uses only the displayed boundary value equation; general spectral existence and completeness are separate analytic results.
\end{definition}

\begin{proposition}[Energy, sign, orthogonality, and scaling]\label{prop:steklov-elementary-identities}
A Steklov eigenfunction satisfies
\[
 \int_\Sigma|\nabla u|^2=\sigma\int_{\partial\Sigma}u^2.
\]
Consequently $\sigma\geq0$; the zero eigenspace consists of constants; and eigenfunctions with distinct eigenvalues are orthogonal in $L^2(\partial\Sigma)$. Under $g\mapsto c^2g$ with $c>0$, every Steklov eigenvalue is divided by $c$.
\end{proposition}
\begin{proof}
Green's identity gives the energy formula. A harmonic function whose boundary trace is zero has zero Dirichlet energy and hence is identically zero; thus the boundary denominator is positive for a nonzero eigenfunction. If $\sigma=0$, zero energy and connectedness make $u$ constant. For eigenfunctions $u,v$ with eigenvalues $\sigma,\tau$, integration by parts in both orders gives
\[
 \sigma\int_{\partial\Sigma}uv=\int_\Sigma\langle\nabla u,\nabla v\rangle
 =\tau\int_{\partial\Sigma}uv.
\]
Finally the Laplacian scales by $c^{-2}$ and the outward unit normal derivative by $c^{-1}$, proving the last assertion directly from the equation.
\end{proof}

\begin{corollary}[A geometric Steklov eigenvalue]\label{cor:free-boundary-steklov-one}
Every nonzero coordinate function of a free-boundary minimal immersion into the unit ball is a Steklov eigenfunction of eigenvalue $1$. Their span has dimension equal to the dimension of the linear span of the image. This assertion alone does not identify the first positive eigenvalue.
\end{corollary}
\begin{proof}
Apply Proposition \ref{prop:free-boundary-coordinate-identities}. A relation $\sum a_iF_i=0$ is exactly a vector orthogonal to the linear span of $F(\Sigma)$, giving the dimension claim. There may be other eigenfunctions with smaller positive eigenvalues; the coordinate calculation supplies no exclusion of them.
\end{proof}

\begin{lemma}[Area acceleration under an ambient flow]\label{lem:free-boundary-flow-acceleration}
Let $\Sigma^n\subset\overline B^N$ be free-boundary minimal and let $X$ be a smooth ambient vector field tangent to the sphere. Put $Y=D_XX$. At $t=0$, the second derivative of area under its flow is
\[
 A''(0)=\int_\Sigma\left[\sum_i|D_{e_i}X|^2+(\operatorname{div}_\Sigma X)^2
 -\frac12\sum_{i,j}\bigl(\langle D_{e_i}X,e_j\rangle+\langle D_{e_j}X,e_i\rangle\bigr)^2\right]
 -\int_{\partial\Sigma}|X|^2,
\]
where $e_1,\ldots,e_n$ is any local orthonormal tangent frame.
\end{lemma}
\begin{proof}
Differentiate the pullback metric along the flow. At an orthonormal frame its first derivative is
$g'_{ij}=\langle D_{e_i}X,e_j\rangle+\langle D_{e_j}X,e_i\rangle$ and
$g''_{ij}=2\langle D_{e_i}X,D_{e_j}X\rangle+\langle D_{e_i}Y,e_j\rangle+\langle D_{e_j}Y,e_i\rangle$.
The second derivative of $\sqrt{\det g}$ is
$\tfrac12\operatorname{tr}g''-\tfrac12\operatorname{tr}((g')^2)+\tfrac14(\operatorname{tr}g')^2$.
After substitution, the additional term is $\int_\Sigma\operatorname{div}_\Sigma Y=\int_{\partial\Sigma}\langle Y,\eta\rangle$ by minimality. On the sphere, differentiate $\langle X,x\rangle=0$ in direction $X$ to obtain $\langle D_XX,x\rangle=-|X|^2$. Since $\eta=x$, this is the claimed boundary term. The flow preserves the sphere and the ball, so the variation is admissible.
\end{proof}

\begin{theorem}[Instability in every positive codimension]\label{thm:free-boundary-ball-instability}
Every nonempty compact free-boundary minimal immersion $\Sigma^n\to\overline B^N$ with $1\leq n<N$ has an admissible ambient variation with strictly negative second variation of area. This is Fraser's instability theorem, \cite{Fraser00}, Theorem 2.9.
\end{theorem}
\begin{proof}
Write the ambient position as $x$ and its tangential projection as $x^T$. For the standard coordinate vectors $E_k$, set
\[
 X_k=2x_kx-(|x|^2+1)E_k,\qquad 1\leq k\leq N.
\]
These fields are tangent to $S^{N-1}$. Their derivatives satisfy
$D_vX_k=2\langle v,E_k\rangle x+2x_kv-2\langle x,v\rangle E_k$.
Thus the symmetric tangential derivative is $4x_kg$, and $\operatorname{div}_\Sigma X_k=2nx_k$. Summing the squared derivative over $k$ and a tangent orthonormal frame gives
\[
 \sum_{k,i}|D_{e_i}X_k|^2=8n|x|^2+4(N-2)|x^T|^2.
\]
On the boundary, $\sum_k|X_k|^2=4(N-1)$. Lemma \ref{lem:free-boundary-flow-acceleration} therefore yields
\[
 \sum_k A''(X_k)=\int_\Sigma[4(N-2)|x^T|^2+4n^2|x|^2]-4(N-1)|\partial\Sigma|.
\]
Because $\Delta|x|^4=8|x^T|^2+4n|x|^2$ and $\partial_\eta|x|^4=4$, integration gives
$|\partial\Sigma|=\int_\Sigma(2|x^T|^2+n|x|^2)$.
Substitution makes the sum equal to
\[
 -4N\int_\Sigma|x^T|^2+4n(n-N+1)\int_\Sigma|x|^2<0.
\]
The second coefficient is nonpositive. At the boundary $x^T=x$ has unit length, so continuity on a boundary collar makes the first integral positive. Hence at least one summand is negative. This calculation follows the conformal fields in \cite{Fraser00}, pp.~968--970, while retaining the exact conormal equality.
\end{proof}

\begin{proposition}[Normal second variation at a free boundary]\label{prop:free-boundary-normal-form}
Let $\Sigma^n\subset\overline B^{n+1}$ be a smooth compact two-sided embedded free-boundary minimal hypersurface, with global unit normal $\nu$. For every smooth scalar function $f$ on $\Sigma$, the second variation in normal direction $f\nu$ is
\[
 Q(f,f)=\int_\Sigma(|\nabla f|^2-|A|^2f^2)-\int_{\partial\Sigma}f^2.
\]
In particular $Q(1,1)<0$. The boundary is allowed to move on the sphere; this quadratic form differs from the fixed-boundary Dirichlet problem.
\end{proposition}
\begin{proof}
The field $f\nu$ is tangent to the sphere along $\partial\Sigma$, because $\eta=x$ is tangent to $\Sigma$. It extends to a smooth ambient vector field tangent to the sphere: use boundary charts in which the transverse intersection of the hypersurface and the sphere is a coordinate intersection, extend its component functions, and combine using a partition of unity. Along $\Sigma$, $D_{e_i}(f\nu)=f_i\nu+fAe_i$. Therefore the squared derivative sum is $|\nabla f|^2+f^2|A|^2$, the tangential divergence is $fH=0$, and the symmetric derivative is $2fh$. Lemma \ref{lem:free-boundary-flow-acceleration} gives the formula. Its value at $f=1$ is strictly negative because $|\partial\Sigma|>0$. Compare \cite{Fraser00}, Section 2.1, p.~960.
\end{proof}

